\chapter{Authentication Layers in zkVM Workflows}

\section{How the Current Encryption Method Was Chosen}

The final method was chosen in two stages. First, primitive benchmarks compared
AES, Salsa20, and LowMC in the same proving stack. Second, the chosen primitive
was integrated into an authenticated pipeline with explicit framing and
verifiable execution.

AES became the practical base because it performed better than LowMC in this
implementation context, and CTR mode offered a straightforward composition path
with HMAC, as described in Chapter~2.
The cipher-selection decision is consolidated at the end of this chapter; the
following sections first explain the authentication layers required once AES-CTR
is used as the confidentiality base.

\section{Threat-Driven Requirements}

The authentication design is driven by three concrete threats:

\begin{itemize}
\item an observer should not learn plaintext from public ciphertext;
\item an adversary should not be able to alter ciphertext, IV, or AAD without tag
verification failing;
\item a verifier should not accept outputs from a different guest program image;
\end{itemize}

The next section maps these threats to separate cryptographic, proof, and
transport layers.

\section{Layer Separation}

This dissertation separates three layers:

\begin{itemize}
\item \textbf{Data authentication:} HMAC over framed data
(\texttt{aad || iv || ciphertext}) detects tampering in the ciphertext path.
\item \textbf{Computation correctness:} zkVM proof verification binds public
outputs to correct execution of a fixed guest image.
\item \textbf{Transport/sender identity:} external signatures or channel
authentication can bind packets to an operational identity.
\end{itemize}

This separation resolves the apparent chicken-and-egg issue. The proof does not
authenticate a network sender because sender identity is not part of the
cryptographic transcript proved by this implementation. That claim would need an
external signature, transport-layer MAC, or authenticated channel. Conversely,
the in-guest MAC does not prove that the zkVM program executed correctly.
Instead, the MAC authenticates the framed data tuple, while the receipt proves
that the expected guest image produced or checked that tuple. The checks are
therefore complementary rather than circular.

In short, a valid proof does not replace message authentication, and a valid MAC
does not prove correct execution of a fixed guest image. Both checks are needed
because the MAC detects transcript tampering, while the receipt binds that
transcript to the fixed guest image. Sender identity remains outside this scoped
model because neither of those checks identifies who transmitted the packet.

\section{Why Encrypt-then-MAC Runs in the VM}

Encrypt-then-MAC runs in the guest rather than as a host-side post-processing
step because both encryption and tag computation are part of the statement being
proved. The verifier should be able to interpret a receipt as evidence that the
fixed guest image encrypted the plaintext, authenticated
\texttt{aad || iv || ciphertext}, and committed the resulting public transcript.
This gives a stronger receipt-binding claim than proving encryption first and
letting the host attach a tag afterwards.

If HMAC were computed only outside the VM, the receipt would no longer bind the
tag to the guest execution trace. A verifier could learn that the guest produced
a ciphertext and that the host later produced a valid MAC, but not that the tag
was produced by the guest under the stated framing rule. Keeping
Encrypt-then-MAC inside the VM makes the ciphertext, IV, AAD, and tag one
authenticated transcript covered by the same receipt.

This placement also clarifies where optional transport authentication belongs. A
deployment may still sign or MAC the final receipt and packet outside the VM to
identify the sender or protect a network channel. That outer authentication is
not a substitute for the in-guest MAC; it answers who sent the packet, while the
in-guest Encrypt-then-MAC path answers whether the ciphertext/tag pair belongs to
the proved computation.

Input preparation, proving orchestration, receipt verification, and benchmark
artifact collection run in the host. Keeping these responsibilities in host code
matches the RISC Zero execution model and keeps deployment concerns separate
from deterministic guest logic \cite{risczeroDocs}.

\section{Alternatives Considered}

\paragraph{AES-GCM.}
AES-GCM combines encryption and authentication in one mode and is widely used,
but this project favored an explicit EtM split because it produced cleaner
component-level overhead attribution and a simpler game-based argument structure
for the dissertation: AES-CTR, HMAC, and the RISC Zero proof each correspond to
one hop in the Chapter 9 IND-CPA game
\cite{nistsp80038d,bellare2000authenticated}.

\paragraph{ChaCha20-Poly1305.}
ChaCha20-Poly1305 is a strong software baseline with broad protocol adoption,
but it was not selected here because the repository and benchmark pipeline were
already centered on AES-derived targets and because introducing a new primitive
family would dilute depth in the final report campaign \cite{rfc8439}.

\paragraph{OCB mode.}
OCB provides efficient one-pass authenticated encryption and is conservative and
standardised --- specified in RFC 7253 with a formal security proof. This project
preferred the explicit EtM split over OCB because the separate CTR and HMAC
components allow direct overhead attribution in the benchmark evaluation and
produce a simpler game-based security argument where each component corresponds
to one hop in the proof \cite{rfc7253,rogaway2001ocb,bellare2000authenticated}.

Table~\ref{tab:auth-alternatives} summarizes the alternatives against the
security and evaluation requirements used for the final selection.

\begin{table}[H]
\centering
\reporttablesetup
\begin{tabularx}{\textwidth}{@{}P{0.16\textwidth}LLL@{}}
\hline
Alternative & Security fit & Evaluation fit & Reason not selected \\
\hline
AES-GCM & Strong AEAD when nonce discipline is correct & Would require
GHASH-specific zkVM measurement & Larger integrated surface and less direct
CTR-vs-auth overhead attribution \\
ChaCha20-Poly1305 & Strong protocol baseline with software-friendly operations &
Would introduce a new primitive family late in the project & Reduces depth of
AES/LowMC comparison and final artifact traceability \\
OCB & Efficient one-pass AEAD with formal support & Not already implemented or
benchmarked in this repository & Additional implementation and licensing/history
discussion not central to research question \\
AES-CTR + HMAC & Standard components with clear EtM reasoning & Directly
supports separate CTR and CTR+HMAC targets & Selected despite two-pass overhead \\
\hline
\end{tabularx}
\caption{Authentication-mode alternatives compared against project requirements.}
\label{tab:auth-alternatives}
\end{table}

\section{Current Implementation Scope}

The implemented path covers the first two layers (HMAC plus zkVM proof
verification). Additional transport-level proof authentication is outside the
current code path.

Therefore, report claims are scoped to confidentiality and integrity of the
cryptographic computation transcript under the fixed-program verification model,
not to sender identity guarantees in a network deployment. They also assume
nonce/IV uniqueness for AES-CTR and do not claim safety if the same counter stream
is reused under the same encryption key.

\section{Final Cipher Selection}

\subsection{Decision \& Justification}

The final cipher selected for the authenticated encryption pipeline is
AES-192 in CTR (Counter) mode for encryption. CTR mode requires only the forward
AES block cipher and unique IV/counter inputs under the same key, as specified in
NIST SP 800-38A and introduced in Chapter~2 \cite{nistsp80038a}.

Authentication is provided by HMAC-SHA256 truncated to 192 bits, using the
standards discussed in Chapter~2 and the in-guest Encrypt-then-MAC design
specified earlier in this chapter. The repository target split used to evaluate
this decision is detailed in Chapter~7.

The decision is based on measured behavior in this zkVM environment. AES
provides better empirical proving performance than LowMC in the tested
implementations, and the AES-CTR path integrates cleanly with the project's
authenticated-pipeline design under FIPS 197 \cite{nistfips197}.

Salsa20 is retained as a benchmarked ARX baseline rather than selected for the
final construction. Its available measurements are older supporting artifacts,
while the final report campaign focused on the AES-versus-LowMC primitive
decision and the matched AES-CTR/CTR+HMAC pipeline.

\subsection{AES Backend Comparison Snapshot}

The remaining engineering question is whether the AES optimization workspace
improves proving performance under strict like-for-like settings.

A direct supporting snapshot from the latest available AES backend comparison
campaign gives:

\begin{table}[H]
\centering
\reporttablesetup
\begin{tabular*}{\textwidth}{@{\extracolsep{\fill}}lrrr@{}}
\hline
Target & Prove s (median) & User cycles (median) & Delta vs \texttt{aes-r0} \\
\hline
\texttt{aes-r0} & 14.8275 & 186256 & baseline \\
\texttt{aes-r0-optimised} & 7.4712 & 96431 & -49.61\% prove, -48.23\% cycles \\
\hline
\end{tabular*}
\caption{AES baseline versus optimisation workspace (single-trial supporting
snapshot from run \texttt{20260401-112048Z}).}
\label{tab:aes-backend-snapshot}
\end{table}

Table~\ref{tab:aes-backend-snapshot} shows a material measured improvement in
that standalone snapshot. Because this comparison is single-trial and outside the
commit-locked final CTR/CTR+HMAC matrix, the final cipher decision is anchored to
the repeated AES-vs-LowMC and authenticated-pipeline results rather than to this
snapshot alone.

The \texttt{aes-r0} value in Table~\ref{tab:aes-backend-snapshot} is therefore
not expected to match the \texttt{aes-r0} value in
Table~\ref{tab:cross-algo-comparison}. Table~\ref{tab:cross-algo-comparison}
uses the May 2026 commit-locked five-trial main report run, whereas
Table~\ref{tab:aes-backend-snapshot} uses the single-trial backend snapshot from
run \texttt{20260401-112048Z}.

A stricter backend study (explicit \texttt{sbox-table} vs \texttt{sbox-logic}
feature flags under repeated trials) remains future work.

\subsection{Decision Summary}

The selected option is not the fastest possible authenticated-encryption design
in general; it is the most defensible design under this project's evidence. AES
beats the available LowMC implementation by a large measured margin, CTR mode
keeps encryption behavior simple to benchmark, and HMAC provides a separate
authentication layer whose overhead can be measured directly. The trade-off is
that the final construction is two-pass and therefore less elegant than an
integrated AEAD mode such as GCM or OCB. That cost is accepted because the
separation improves artifact traceability and simplifies the security argument.
